\documentclass[10pt,a4paper]{amsart}
\usepackage[margin=19mm]{geometry}
\usepackage{amsmath,amssymb,mathtools}
\usepackage[scr=rsfs,cal=euler]{mathalfa}
\usepackage{xfrac}
\usepackage[unicode,hidelinks,bookmarks=false]{hyperref}
\newcommand{\ie}{i.e.\ }
\newcommand{\cf}{cf.\ }
\newcommand{\resp}{respectively\ }
\newcommand{\Subsection}{\subsection}
\providecommand{\nobreakdash}{\nobreak\hbox{-}}
\setlength{\emergencystretch}{2em}
\multlinegap0pt
\usepackage{bm}
\usepackage{bbm}
\usepackage{dsfont}
\usepackage{xspace}
\usepackage{cancel}
\usepackage{mathtools}
\usepackage{amsthm}
\makeatletter
\@ifundefined{lemma}{\newtheorem{lemma}{Lemma}}{}
\makeatother
\DeclareMathOperator{\E}{\mathbb{E}} % Expectation operator
\title[Toward Optimal ANC: Establishing Mutual Information Lower Bo]{Toward Optimal ANC: Establishing Mutual Information Lower Bound}
\author{Fran\c{c}ois Derrida, Shahar Lutati, Eliya Nachmani}
\date{Source: arXiv:2505.17877v1 (CC BY 4.0)}
\begin{document}
\maketitle

\noindent\emph{Source and licence.} Toward Optimal ANC: Establishing Mutual Information Lower Bound. Version arXiv:2505.17877v1 (CC BY 4.0), distributed under the Creative Commons Attribution 4.0 International licence, as recorded in the arXiv licence metadata for this exact version and shown on the abstract page https://arxiv.org/abs/2505.17877v1. Full rights notice: licenses/arxiv-2505.17877-CC-BY-4.0.txt.\par\smallskip

\noindent\textit{Editorial scope.} Counted unit: the Lemma whose source label is \texttt{eq:info\_bound\_raw} in the article (source0.tex lines 164--171, source label \texttt{eq:info\_bound\_raw}), delivered together with the article's own complete original proof of it (source0.tex lines 173--224, one \texttt{proof} environment of 52 lines). No mathematics is written, repaired or re-derived: statement and proof are the article's own text line for line. Required context retained uncounted: none. Every definition, display and label of those lines is retained unchanged. Omitted from this package and not redistributed: 1--163, 172--172, 225--449. Nothing inside a retained excerpt is omitted or altered: every retained body is delivered exactly as the article's lines are written, so no omission receipt of any kind accompanies this package. Citations: the retained excerpts carry no citation command at all, so no bibliography entry of the article is transcribed and no reference is redistributed. Reference adaptation: none was needed and none was made; every reference inside the delivered unit resolves inside it, so no unresolved reference or citation is produced. Printed slips kept literal: no printed word, symbol, hypothesis or number of the counted statement or of its proof was altered. Layout and class: the article's own class is \texttt{article} with packages including neurips\_2025, amsmath, inputenc, fontenc, hyperref, url, booktabs, amsfonts, nicefrac, microtype, xcolor, graphicx, subcaption, wrapfig; the wrapper is the lane's pinned \texttt{amsart} editorial wrapper at 10pt on A4, which restates the theorem-like environments the retained text needs and adds the article's own macro definitions that the retained text needs (\texttt{\textbackslash E}), with extra wrapper packages (bm, bbm, dsfont, xspace, cancel, mathtools). The delivered numbering is the unit's own. Assets: no retained span contains a figure environment, a graphics-inclusion command, a PDF-page inclusion, a table or a TikZ picture; no image file is copied into this package, the package declares no asset dependency and no image file is redistributed with it. Licence: version arXiv:2505.17877v1 of this article is distributed under the Creative Commons Attribution 4.0 International licence, as recorded in the arXiv licence metadata for this exact version and shown on the abstract page https://arxiv.org/abs/2505.17877v1. A full rights notice is delivered at licenses/arxiv-2505.17877-CC-BY-4.0.txt. Characters: every retained excerpt is pure ASCII, so the delivered engine, XeLaTeX with its default Latin Modern text font, reports no missing character.
\begin{lemma}[Information-Theoretic Bound on Cancellation Error]
Under the assumption that \( y(n) \) and \( d(n) \) are jointly WSS and ergodic, and that \( d(n) \) possesses a finite differential entropy rate \( H(d) \), the average power of the residual error \( e(n) = d(n) - f(y(n)) \) is lower bounded by:
\begin{equation}
\sigma_e^2 = \E[|e(n)|^2] \geq \sigma_d^2 \left( 1 - \frac{I(y(n); d(n))}{H(d(n))} \right)
\label{eq:info_bound_raw}
\end{equation}
where \( \sigma_d^2 = \E[|d(n)|^2] \) is the average power of the disturbance signal.
\end{lemma}

\begin{proof}
Let \( d(n) \) and \( y(n) \) be jointly wide-sense stationary (WSS) and ergodic random processes. Suppose \( d(n) \) has finite variance \( \sigma_d^2 \) and finite differential entropy rate \( H(d) \). Define the residual error as:
\[
e(n) = d(n) - f(y(n))
\]
where \( f: \mathbb{R} \to \mathbb{R} \) is any deterministic mapping (possibly nonlinear). Let \( \hat{d}(n) = f(y(n)) \) be the estimate of \( d(n) \) based on \( y(n) \).

From information theory, the mutual information between \( d(n) \) and \( y(n) \) satisfies:
\begin{equation}
I(d(n); y(n)) \geq R(D)
\label{eq:rate_distortion_lower_bound}
\end{equation}
where \( R(D) \) is the rate-distortion function of the source \( d(n) \) at distortion level \( D = \mathbb{E}[(d(n) - \hat{d}(n))^2] = \sigma_e^2 \).

For general stationary ergodic sources with finite differential entropy rate \( H(d) \), the Shannon lower bound on the rate-distortion function provides:
\begin{equation}
R(D) \geq H(d) - \frac{1}{2} \log(2\pi e D)
\label{eq:shannon_lower_bound}
\end{equation}
Solving \eqref{eq:shannon_lower_bound} for \( D \) yields:
\begin{equation}
    D \geq \frac{1}{2\pi e} \cdot \exp\left(2(H(d) - I(d(n); y(n)))\right)
    \label{eq:inequality_exp}
\end{equation}

However, this expression is difficult to interpret directly. Instead, we derive a more interpretable but looser bound by noting that mutual information is always upper bounded by entropy:
\[
I(d(n); y(n)) \leq H(d)
\]
Hence, we can define the normalized mutual information ratio:
\[
\alpha := \frac{I(d(n); y(n))}{H(d)} \in [0, 1]
\]
Recall, that for gaussian noise, the entropy is maximazied and the following holds:
\[
\sigma^2_d  = \frac{1}{2\pi e}e^{2H(d)} \]
Rearranging terms, of \eqref{eq:inequality_exp}, and keep in mind that the ratio is always smaller than one, so for negative exponent it is upper bound,
\[
D \geq \frac{1}{2\pi e} \cdot \exp\left(2H(d)\right) \cdot \exp\left(- \frac{I(d(n); y(n))}{H(d)}\right)
\]
expanding the second exponent as taylor series, and taking the first order yeilds,
\begin{equation}
\sigma_e^2 \geq \sigma_d^2 \left(1 - \frac{I(d(n); y(n))}{H(d)} \right)
\end{equation}
Intuitively, this ratio measures the fraction of the information in \( d(n) \) that is captured by \( y(n) \). If \( \alpha = 1 \), perfect reconstruction is possible in principle, and \( \sigma_e^2 \to 0 \). If \( \alpha = 0 \), then \( y(n) \) carries no information about \( d(n) \), and \( \sigma_e^2 \geq \sigma_d^2 \).

We thus propose a linear lower bound on the distortion based on this ratio:
\begin{equation}
\sigma_e^2 \geq \sigma_d^2 (1 - \alpha) = \sigma_d^2 \left(1 - \frac{I(d(n); y(n))}{H(d)}\right)
\end{equation}
which completes the proof.
\end{proof}

\end{document}
